\silentchapter{Fundament}{fundament}
\begin{widebody}
{\small

Rammeverket er ein heurestikk. Det seier ikkje kva forma blir; det seier kva ein må sjå etter. Slik er all operasjonell vitskap: pavlovsk læring er ein heurestikk som let seg operasjonalisere på kvart nivå frå genregulatoriske nettverk til sivilisasjonar, og verdien ligg i at han gjer det same spørsmålet handterleg på tvers av skalaer. Fri-energi-prinsippet\textsuperscript{19,\,83} gjer det same for agens. Prinsippet er per 2026 empirisk demonstrert på tvers av minst fjorten substratnivå. Eit utdrag syner spennvidda.

\bigskip
\begin{center}
\footnotesize
\renewcommand{\arraystretch}{1.2}
\begin{tabular}{@{}p{38mm}p{50mm}l@{}}
\toprule
\textbf{Skala} & \textbf{Mekanisme} & \textbf{Kjelde} \\
\midrule
Genregulatorisk nettverk & Seks pavlovske læringsformer & \textsuperscript{36} \\
Einsleg celle & Habituering, assosiativ betinging & \textsuperscript{98} \\
Slimsopp & Habituering, minneoverføring & \textsuperscript{99} \\
Vev & Bioelektrisk morfogenetisk minne & \textsuperscript{85} \\
Immunsystem & Trena immunitet via epigenetikk & \textsuperscript{100} \\
Hjerne & Synaptisk plastisitet, prediktiv prosessering & \textsuperscript{101} \\
Individ & Seks pavlovske former, symbolsk læring & \textsuperscript{102} \\
Dyade & Bayesiansk synkroni & \textsuperscript{103} \\
Institusjon & Utforsking og utnytting & \textsuperscript{104} \\
Sivilisasjon & Paradigmeskift som kollektiv nylæring & \textsuperscript{105} \\
\bottomrule
\end{tabular}
\end{center}
\bigskip

At den same formalismen prediserer fenomen over ein skala frå nanometer til kontinent er ei urimeleg effektivitet i Wigners\textsuperscript{106} forstand. Ho rettferdiggjer pragmatisk bruk sjølv før dei djupaste tolkingane er avgjorte.

Ein observasjon om kva slags framgang dette er. Vitskapshistoria det siste hundreåret kan lesast som ein systematisk oppløysing av binære kategoriar. Materie og energi vart éi storleik; bølgje og partikkel vart dualitet; gen og miljø vart epigenetikk; hjerne og kropp vart ein kroppsleggjort kognisjon; medvit og ikkje-medvit vart gradert agens; levande og ikkje-levande vart autopoietisk kontinuum. Kvar gong ein hard dikotomi vart rekna ut som ein heurestikk heller enn ein ontologi, opna det eit forskingsrom som hadde vore lukka. Formlære står i denne linja. Form og funksjon, natur og artefakt, intensjon og emergens, biologisk og kulturell, menneskeleg og teknologisk, formgjevar og brukar: alle er heurestiske skilje, ikkje kategoriske. Å handsame dei som gradientar i staden for grenser er ikkje berre å låne eit grep frå naturvitskapen. Det er å følgje den einaste rørsla som har produsert varig framgang.

Designhandlinga er formelt isomorf med aktiv inferens: ein agent med ein generativ modell av ei form som skal bli, ei observasjon av den forma som er, og ein justeringsmekanisme mellom dei. Formalismen er substrat-uavhengig. Det designfaget hittil har mangla, finst ferdig formulert i nabofaga. Denne seksjonen tek ikkje opp noko nytt; han peikar på det som ligg og ventar.

Dette gjev faget eit forskingsprogram med fire punkt. \textit{Fyrst}, å definere kjerneomgrepa operasjonelt, med referanse til eit rammeverk som alt er validert. Form, formgjevar og designar er ikkje lenger fritt definerte; dei er posisjonar i eit tilpassingslandskap, agentar som navigerer det, og arter av agent med spesifikke måltilstandar. \textit{Andre}, å operasjonalisere omhug som mengda av aksar ein agent aktivt justerer for, og utvikle kvantitative mål; K-metrikken\textsuperscript{86} gjev éin slik. \textit{Tredje}, å ta i bruk dei enorme empiriske arkiva museum, bibliotek, etnografar og statistikarar har samla over generasjonar. Desse arkiva er ikkje mangelvare; dei er underutnytta, ofte tolka med kategorifeil der stil vert handsama som årsak, funksjon som primærtrykk, og brukaren som den einaste agenten. Reanalysering av eksisterande designforskingsdata under denne lins er eit konkret forskingsprogram. \textit{Fjerde}, å binde fagleg sertifisering til dokumentert omhugsbreidde, slik medisin bind sertifisering til dokumentert verknad.

Materialet ligg der. Formlæras bidrag er å binde det saman.

}
\end{widebody}
